\section{Introduction}

This paper is concerned with the efficiency of matching in decentralized matching markets.
In particular, many markets with horizontal differentiation nevertheless lack adequate mechanisms for participants to signal heterogeneity of preferences to the other side, leaving them to exert costly effort to overcome the resulting search friction.
This can lead to inefficient matching: participants do not acquire sufficient information to find their ``best match'' on the other side.

In response to this challenge, markets might consider the introduction of signaling mechanisms that allow one side of the market to reveal preference information to the other side.
Indeed, economists have long-recognized the potential for signaling to overcome market informational problems \citep{spence1973job, crawford1982strategic}.
However, it is far from clear that the introduction of such a mechanism would improve outcomes.  There are at least two salient concerns.
First, parties may not even choose to use the possibility to credibly signal; we refer to this outcome as the {\em pooling} equilibrium, in the sense that the signal does not distinguish agents.  In this case, of course, the introduction of a signal clearly has no welfare effect.
Second, even if a {\em separating} equilibrium is obtained where participants credibly signal, the participants are strategic; therefore, there is no guarantee that the separating equilibrium will be more efficient than the previous pooling equilibrium.  

In this paper, we investigate these issues by studying a randomized controlled experiment of a new signaling feature in a large online labor market \citep{Horton2010}.  In the setting we consider, applicants are vertically differentiated in experience, ability, expertise, etc.; however, employers have heterogeneous preferences over the relative importance of prices (wages) and these quality dimensions.  
The signaling mechanism introduced into the market gave employers a means to publicly state their preferences to potential applicants, revealing their relative willingness to pay for ``better'' workers.
In principle, such a mechanism could allow employers who are willing to pay more for the most qualified workers to signal as much to those applicants.
Of course, by signaling in this way, employers risk that applicants' strategic response will be to take advantage of the information through increased wage bids.

The signaling feature was simple: when posting a job, employers had to select from one of three ``tiers'' to describe the kinds of applicants they were most interested in:
(1) Entry level: ``I am looking for workers with the lowest rates.'';
(2) Intermediate: ``I am looking for a mix of experience and value.'';
(3) Expert: ``I am willing to pay higher rates of the most experienced workers.''
We refer to these tiers are ``low,'' ``medium,'' and ``high,'' respectively.
When introduced market-wide, this signal would be revealed publicly to all job-seeking workers. 
However, during the experimental phase, not all firms had their signals revealed.
By selectively revealing some signals, we could learn more about the effects of revelation.\footnote{
  We will use the terms ``firm'' and ``employer'' interchangeably. Also, our use of the word employer simple reflects the typical usage of these terms in the literature and is not a comment on the nature of the legal relationship created on the platform. 
} 

Observe that if employers that signal ``high'' are indicating that they are willing to pay for the ``best'', most proven workers; by contrast, employers that signal ``low'' are indicating that they are relatively more sensitive to price, and thus are willing to hire unproven or less able (and ostensibly lower price) workers.  
Firms must consider the effects of truthfully revealing their preferences, since such signaling will affect both (1) the pool of applicants they receive; and (2) the wage bids of those applicants. 
Since workers can condition their wage bids on the preferences declared by employers, it is not clear that truthful revelation will be in the employer's best interest.
In particular, a firm declaring ``high'' may get higher quality applicants, but at the cost of higher wage bids by those applicants.
Conversely, a firm declaring ``low'' might benefit from lower wage bids, but at the cost of a lower quality applicant pool.

Our main contributions revolve around addressing the central issues raised above: Does a separating equilibrium obtain upon introduction of the signal?
And if such an equilibrium exists, what are the consequences for welfare?  
First, we provide evidence that in fact a separating equilibrium is obtained: employers have strong heterogeneity in their preferences; they declare these preferences through use of the signal; and applicants sort and adjust their wage bids in response to employers' declarations.
Second, we provide evidence that this equilibrium is in fact welfare improving, by studying improvement in the quality of match outcomes in the new separating equilibrium relative to the prior pooling equilibrium.

Our data comes primarily from an experimental period, in which employers were assigned to different treatment groups.
During the experimental phase, employers were randomized to one of two  arms,  each of which was split into two experimental groups. 
In the ``explicit'' arm of the experiment, employers were told, \emph{ex ante}, that their signal would either be \emph{shown} or \emph{not shown} to applying workers. 
In the ``ambiguous'' arm of the experiment, employers were told the platform \emph{might} reveal their preferences to workers, and then employers were randomized to either have their preferences revealed to would-be job applicants or not, \emph{ex post}.
Any subsequent job opening posted by an employer had the same treatment assignment as the original opening.
However, unless otherwise noted, we only use the first opening posted by the employer. 
We also make limited use of some post-experimental data, with our goal being to assess whether a separating equilibrium persisted once all signals were revealed. 

We organize our research findings as follows.  First, we discuss heterogeneity of employer preferences, and their signaling behavior via the feature.  Next, we show that applicants sort and adjust their wage bids in response to this preference revelation, completing our confirmation that separation obtains upon introduction of the signaling mechanism.  Finally, we study match outcomes and the welfare implications of a separating equilibrium.

\begin{enumerate}
\item {\bf Employer preferences and signaling behavior.}  We begin by using the explicit arm of the experiment, where employers were told their preferences would not be shown to applicants; this takes away any impact of strategic considerations on preference revelation.
Studying the behavior of these employers reveals that there is substantial heterogeneity in price/quality preferences.
Some of this variation can be explained by differences in preferences across categories; for example, in the ``Administrative Support'' category, \fracLowAdminSpt{}\% employers chose the low tier signal compared, to just \fracLowWeb{}\% in ``Web Development'' making the same choice; in ``Networking \& Information Systems'' the fraction of employers choosing the high tier is \fracHighNetworking{}\%, whereas it is only \fracHighAdminSpt{}\% in ``Administrative Support.''
However, despite the evident importance of the job category, there is still substantial within-category variation, with no potential tier signal garnering even close to 100\% of employers.
The residual within-category variation is important, as it implies that the signal \emph{could be} informative to workers and not simply subsumed by other sources of information.

We then use both the explicit and ambiguous arms to demonstrate that employers {\em do} in fact declare these preferences via the signaling mechanism.  
Indeed, we find that preference revelation by employers nearly exactly 
Further, their willingness to reveal preferences is not affected by whether or not they are told applicants will be shown their preferences.

Illustrating this cross-category variation, in the ``Administrative Support'' category, \fracLowAdminSpt{}\% employers chose the low tier signal compared, to just \fracLowWeb{}\% in ``Web Development'' making the same choice; in ``Networking \& Information Systems'' the fraction of employers choosing the high tier is \fracHighNetworking{}\%, whereas it is only \fracHighAdminSpt{}\% in ``Administrative Support.''
Despite the evident importance of the job category, there is still substantial within-category variation, with no potential tier signal garnering even close to 100\% of employers.
The residual within-category variation is important, as it implies that the signal \emph{could be} informative to workers and not simply subsumed by other sources of information.

\item {\bf Applicant response to employer preference revelation.}  Next, we investigate applicants' response to employer preference revelation; in particular, whether or not they sort by relative quality, and how they adjust their wage bids.
For this analysis, we use the  ``ambiguous arm'' of the experiment, in which employers were told that the platform \emph{might} reveal their preferences to workers, with actual revelation determined at random, \emph{ex post}. 
Because employers did not know \emph{ex ante} whether their signal would be revealed, we can treat their signal as exogenous, allowing us to see how revelation of the signal affected the applicant pool, holding the tier selection fixed.
Using such an approach we find that workers strongly respond to the signals sent by employers: in particular, employers that signal a higher preference for quality vs. price do indeed receive a relatively more experienced applicant pool.
(We do not find significant evidence that the size of applicant pools is different in the different tiers.)


Regarding wage bids, we find that workers strongly condition their wage bids on the employer's revealed signal.
For high tier employers, revelation of the signal increased the wage bids they received by 10\%.
For medium tier employers, there was no detectable effect, while for low tier employers, revelation lowered wage bids by 15\%.
We further show that this change in bids is not solely compositional. 
In particular, using the fact that workers send multiple applications, we show that \emph{within} workers, there is direct conditioning on the revealed signal---the same worker bids more when they know they are facing a employer with ``higher'' vertical preferences.
%Furthermore, we present evidence that shading in wage bids was not due to perceived job amenities/disamenties that might correlated with the employer tier, but rather that workers were considering employer willingness to pay directly.

\item {\bf Match outcomes}.  Finally, we investigate  match outcomes.  First, while we find little change in the level of hiring, there are large effects on the composition of hired workers. 
The main composition effect was for substantially less experienced workers to be hired by low tier employers:
hired workers had about 12\% less experience and about 18\% lower past hourly earnings when the signal was revealed.
There is little compositional change from revelation in the other tiers, at least with respect to observable worker attributes.
In terms of the wages of hired workers, in both the medium and high tiers, wages were 6\% higher following revelation, and about 10\% lower in the low tier. 

These changes by themselves leave ambiguous whether match outcomes were improved, and in particular, whether the introduction of the change is welfare improving.  
If we take the simplified view that employer surplus is the difference of their utility for work completed less payment, and that applicants have no hidden costs, then aggregate welfare is simply the aggregate utility of employers.
We focus on two outcomes: total hours worked, and the total wage bill.
Increases in these quantities would suggest that the matching of employers to workers was welfare improving: employers can choose to terminate a contract at any time, so the willingness to engage in a longer contract at higher expense suggests greater employer satisfaction with the matches made.
 
We find that overall, hours-worked increased, as did the total wage bill.
Hours-worked increased by about 5\% and the total wage bill by about 3\%.
These increases occurred in every tier. 
An increase in hours-worked is unsurprising in the low-tier, given the lower wages. 
However, an increase in hours-worked in the high-tier is surprising: 
hired workers in the high tier had nearly 5\% higher wages despite not being \emph{observably} more productive, and yet we observe more hours-worked.
At higher wages, firms should economize on labor, and yet if anything, they demanded more hours.
A reconciliation is that revelation lead to better matches, and so despite the higher wage, employers demanded more hours.

One counter-argument to this ``more hours, better matches'' argument is that with higher wages, workers have an incentive work more hours, stretch out tasks and pad hours. 
However, this seems unlikely, as we also find that employers who had their signal revealed---and hence were paying higher wages but also buying more hours---gave higher feedback to the worker and to the platform.
Furthermore, on the platform, employers have extensive tools to monitor and cap hours-worked, suggesting we should think of employers as choosing the number of hours. 

To the extent these measures are reasonable proxies for employer surplus, it seems likely that match quality improved. 
Although there are no theoretical guarantees that a separating equilibrium is preferable, in this platform, the results are consistent with the platform goals:
employers received more of the right kinds of applicants and seemed to form better matches. 
\end{enumerate}

In terms of generality, the kind of problem this signaling mechanism addresses is commonplace in markets.
Relatively few markets are characterized by fully-formed preferences and centralized match-making.
This can be seen in the growing interest in how preferences are learned in matching markets \citep{acemoglu2017, ashlagi2017communication}.
Many more markets require buyers and sellers seek each other out, initiating contact with potential partners and then deciding whether to match in a decentralized fashion.
In these markets, both sides engage in costly search and must weigh the attributes of the other side against the proposed terms---a key market design challenge is lowering these costs.
As more of economic life becomes computer-mediated \citep{varian2010computer}, introducing signaling mechanisms is likely to be an important tool in the market designer ``toolbox.''
The ``design pattern'' illustrated in this paper can be applied to many kinds of marketplaces---platforms can take some dimension on which buyers and sellers differ and attributes and then offering the two sides structured ways to send and receive such signals. 

The paper is organized as follows.
In Section \ref{sec:empirical_context}, we provide context for the online labor market setting of the experiment we study.
In Section \ref{sec:exp_design}, we outline the experimental design.
[[[ COMPLETE THIS... ]]]
